\section{Metodología}
\label{sec:metodologia}

\subsection{Dise\~no general}

El estudio adopta un dise\~no observacional, retrospectivo y no experimental. No se
intervino sobre la comunidad ni se modificó su comportamiento; se analizó, en cambio,
un corpus ya existente y generado de forma natural. El enfoque es \emph{exploratorio y
descriptivo}: su propósito es caracterizar el corpus y cuantificar patrones, no probar
hipótesis causales.

La unidad de análisis es el mensaje. Para cada mensaje se calculan dos salidas: una
etiqueta de polaridad (positivo, negativo, neutro o indeterminado) y un vector de
activación afectiva multidimensional. Adicionalmente, cada mensaje recibe
una marca de nivel de evidencia que documenta la cantidad de señal disponible.

\subsection{Arquitectura del \emph{pipeline}}

El \emph{pipeline} base se implementa en R \citep{rcore2024} y se compone de diez scripts
encadenados, cada uno de los cuales consume los artefactos del anterior y produce nuevos
artefactos. La Tabla~\ref{tab:etapas} resume las etapas.

\begin{table}[htbp]
  \centering
  \small
  \caption{Etapas del \emph{pipeline} de análisis.}
  \label{tab:etapas}
  \begin{tabular}{clp{8.2cm}}
    \toprule
    \# & Etapa & Descripción \\
    \midrule
    01 & Importación & Lee el JSON de Telegram, filtra servicio y multimedia sin texto, persiste mensajes. \\
    02 & Validación & Verifica integridad, unicidad y coherencia temporal; marca duplicados de contenido. \\
    03 & Limpieza & Normaliza el texto (Unicode NFC, minúsculas, espacios) conservando los emojis. \\
    04 & Sentimiento & Procesa emojis y negación; calcula polaridad, emociones y nivel de evidencia. \\
    05 & Tópicos & Estima el modelo LDA sobre la matriz documento--término. \\
    06 & Temporal & Agrega por día, semana y mes; series de polaridad y volumen. \\
    07 & Markov & Estima la cadena de Markov sobre la secuencia de estados. \\
    08 & Evaluación & Calcula cobertura y, si existen, métricas contra etiquetas manuales. \\
    09 & Informe & Genera figuras, tablas, macros y compila el PDF. \\
    10 & Pruebas & Ejecuta la batería de pruebas automatizadas. \\
    \bottomrule
  \end{tabular}
\end{table}

Cada etapa abre una transacción en la base de datos de auditoría SQLite y registra, al
concluir, el número de elementos procesados, la marca de tiempo y los artefactos
producidos. Esto garantiza la trazabilidad completa: es posible reconstruir qué versión
del código y de los datos produjo cada figura del informe. 

Para robustecer la etapa 04 y dotarla de capacidad operativa en entornos de producción, se integró el motor modular \textit{Emotica}, cuyo flujo detallado de inferencia se esquematiza en la Figura~\ref{fig:pipeline_emotica}.

\input{figures/fig_pipeline}

\subsection{Extracción y transformación de datos}

La exportación de Telegram Desktop produce un archivo JSON con dos tipos de entradas:
mensajes de usuario y mensajes de servicio (avisos de ingreso, cambios de título, etc.).
La etapa de importación distingue ambos, descarta los segundos y los mensajes
multimedia sin texto, y conserva el resto. El texto puede aparecer como una cadena de
caracteres o como una lista de segmentos heterogéneos (texto plano, enlaces, menciones y
emojis); el importador normaliza ambas formas a una representación textual uniforme y,
por separado, extrae la lista de emojis y el diccionario de reacciones.

Las marcas de tiempo se almacenan en UTC y se conservan también en hora local. La
normalización del texto aplica la forma Unicode NFC, convierte a minúsculas y colapsa
espacios en blanco, con la decisión de dise\~no crucial de \emph{no} eliminar los emojis,
que permanecen como unidades léxicas de pleno derecho. Asimismo, se incorpora un paso previo de anonimización PII (\textit{Personally Identifiable Information}) mediante el módulo \texttt{pii.R}, sustituyendo números de teléfono, correos electrónicos y menciones por tokens neutros (\texttt{[PHONE]}, \texttt{[EMAIL]}).

\subsection{Se\~nales de evidencia}

El clasificador separa tres fuentes de evidencia, resumidas en la Tabla~\ref{tab:senales}.

\begin{table}[htbp]
  \centering
  \small
  \caption{Se\~nales de evidencia y su papel en el clasificador.}
  \label{tab:senales}
  \begin{tabular}{lp{5.5cm}p{5.5cm}}
    \toprule
    Se\~nal & Origen & Tratamiento \\
    \midrule
    Léxico & Palabras y frases del diccionario & Peso y polaridad por término, con inversión por negación. \\
    Emojis & Caracteres emoji del texto & Entrada léxica propia con polaridad y emoción; peso configurable ($\alpha$). \\
    Reacciones & Contadores de reacción de Telegram & Se registran; peso por defecto nulo para evitar doble conteo. \\
    \bottomrule
  \end{tabular}
\end{table}

\subsection{Puntuación de polaridad}

La puntuación de polaridad de un mensaje se define como la suma algebraica de los pesos
de los términos reconocidos, más la contribución de los emojis ponderada por un
coeficiente $\alpha$:

\[
\text{score}(m) \;=\; \sum_{i \in \text{l\'exico}(m)} w_i \,\cdot\, s_i
   \;+\; \alpha \sum_{j \in \text{emojis}(m)} w_j ,
\]

\noindent donde $w_i$ es la magnitud del término, $s_i \in \{-1,+1\}$ es su signo
(invertido si cae en la ventana de negación) y $w_j$ es la polaridad del emoji. La
etiqueta de polaridad resulta del signo del \emph{score}: positivo, negativo o neutro.

\subsection{El motor afectivo Emotica: taxonomía de 15 emociones y din\'amica de opuestas}
\label{sec:emotica_motor}

Para trascender la polaridad unidimensional y mapear estados afectivos complejos, el motor \textit{Emotica} implementa una taxonomía de quince emociones discretas organizadas en dos dimensiones psicológicas: \textbf{Valencia} (positiva, negativa o ambigua) y \textbf{Activación o Arousal} (baja, media o alta), declaradas en \texttt{resources/taxonomy/emotions.yml}. La Tabla~\ref{tab:taxonomia_15} y la Figura~\ref{fig:espacio_afectivo} detallan esta distribución.

\input{tables/tab_taxonomia_15}

\input{figures/fig_espacio_afectivo}

Un componente fundamental del motor es la \textbf{transferencia hacia emociones opuestas}. Cuando un término asociado a una emoción $e$ aparece precedido por un negador dentro de una ventana contextual de tres tokens ($w=3$), el sistema transfiere una fracción del peso a la emoción opuesta configurada ($e_{op}$):
\begin{equation}
    S(e_{op}) \leftarrow S(e_{op}) + \lambda_{neg} \cdot S(e, t)
    \label{eq:opuesta}
\end{equation}
donde $\lambda_{neg} = 0.5$ (\texttt{factor\_negacion\_opuesta}). De este modo, una locución como \textit{"no estoy tranquila"} no se clasifica como neutra ni vacía, sino que activa directamente la categoría de \textbf{Preocupación}.

\subsection{Saturación no lineal y se\~nales estilísticas}
\label{sec:saturacion}

A fin de prevenir que mensajes extensos o repetitivos distorsionen la escala afectiva, Emotica sustituye la acumulación puramente lineal por una función de saturación exponencial normalizada en el intervalo $[0, 1]$:
\begin{equation}
    I(e) = 1 - \exp\left(-\frac{S(e)}{\sigma}\right) + \Delta_{estilo}
    \label{eq:saturacion_exp}
\end{equation}
donde $\sigma = 2.0$ (\texttt{escala\_intensidad}) y $\Delta_{estilo} \le 0.60$ es un factor de intensificación tipográfica configurado en \texttt{config/engine.yml} (Tabla~\ref{tab:reglas_motor}), el cual premia la presencia de signos de exclamación (+0.15 hasta un tope de 3), palabras completas en mayúsculas (+0.20) y alargamiento vocálico informal (+0.15, ej. \textit{"graciaaaas"}).

\input{tables/tab_reglas_motor}

El valor final $I(e)$ se categoriza en los niveles discretos \textit{Leve} ($[0.30, 0.60)$), \textit{Moderada} ($[0.60, 0.85)$) y \textit{Alta / Muy Alta} ($\ge 0.85$).

\subsection{Detección multipalabra de jerga regional e idioms}
\label{sec:jerga}

En el discurso universitario e informal de mensajería, gran parte de la carga emocional radica en locuciones compuestas (\textit{n-grams}) que pierden su semántica si se tokenizan como palabras aisladas. Emotica introduce una etapa de detección previa de modismos mediante los archivos \texttt{jerga\_modismos.csv} y \texttt{slang\_idioms.csv}. 

El motor cubre expresiones locales del contexto universitario boliviano (UMSA) como \textit{"zafé el semestre"} (alegría/tranquilidad), \textit{"me jalaron"} (tristeza/decepción), \textit{"barco"} o \textit{"reventado"}, junto con modismos hispanoamericanos (México, España, Argentina, Chile, Colombia) y abreviaturas de chat (\textit{tqm, ntp, qepd}). Asimismo, se contemplan locuciones idiomáticas en inglés como \textit{"over the moon"}, \textit{"last straw"} o \textit{"butterflies in my stomach"}.

\subsection{Extensión multilingüe y filtrado de ruido}
\label{sec:multilingue}

La arquitectura de Emotica es estrictamente modular y desacoplada del código ejecutable. Para habilitar un nuevo idioma basta con suministrar su carpeta correspondiente en \texttt{resources/}, sin alterar el núcleo del algoritmo. La Tabla~\ref{tab:cobertura_lexica} resume el inventario léxico activo en español e inglés.

\input{tables/tab_cobertura_lexica}

Adicionalmente, se integra un catálogo universal de 75 emojis afectivos (\texttt{\_universal/emojis.csv}) válido para cualquier lengua, y filtros heurísticos de publicidad (\texttt{resources/anuncios/\{es,en\}.txt}) que descartan mensajes comerciales y convocatorias masivas para evitar falsos positivos en emociones como alegría o esperanza.

\subsection{Nivel de evidencia y la etiqueta \emph{indeterminado}}

El n\'umero de pruebas $e(m)$ es el total de términos y emojis reconocidos en el
mensaje. Se define:

\[
\text{evidencia}(m) =
\begin{cases}
  \text{sin\_evidencia}, & e(m) = 0,\\
  \text{insuficiente},   & 0 < e(m) < \tau,\\
  \text{suficiente},     & e(m) \ge \tau,
\end{cases}
\]

\noindent con umbral $\tau$ configurable (por defecto $\tau = 1$). Si un mensaje no
aporta evidencia, su etiqueta de polaridad se fuerza a \emph{indeterminado} y nunca a
\emph{neutro}. Esta regla expresa un principio metodológico: la ausencia de evidencia no
es evidencia de ausencia, y la neutralidad debe reservarse para mensajes con se\~nal
explícita cuyo balance es cercano a cero.

\subsection{Análisis temporal}

La etapa temporal une los mensajes con sus resultados de sentimiento, reconstruye las
marcas temporales y agrega por día, semana y mes. Para cada período se calculan el
volumen, la polaridad media, los conteos por etiqueta, las proporciones normalizadas y
los totales por emoción. Se aplican además medias móviles de suavizado para distinguir
tendencia y ruido.

\subsection{Modelado de tópicos}

Los tópicos se estiman con LDA \citep{blei2003lda} mediante \texttt{quanteda}
\citep{benoit2018quanteda} y \texttt{topicmodels} \citep{grun2011topicmodels}. El
preprocesamiento construye una matriz documento--término tras eliminar puntuación,
n\'umeros y palabras vacías del espa\~nol, con n-gramas de uno a dos términos. El
n\'umero de tópicos $k$ se fija en la configuración y el muestreo usa una semilla fija
para garantizar reproducibilidad. Los resultados se interpretan cualitativamente, tal
como recomienda \citet{grimmer2013text}.

\subsection{Cadenas de Markov}

La secuencia de estados de polaridad se modela como una cadena de Markov de primer orden
\citep{norris1997markov}. La matriz de transición $P$ se estima por máxima verosimilitud
a partir de las transiciones observadas entre mensajes consecutivos. A partir de $P$ se
calculan la distribución estacionaria $\pi$ y el tiempo medio de permanencia en cada
estado. Estos indicadores describen la dinámica: un estado con alta permanencia act\'ua
como un atractor del sistema.

\subsection{Marco de evaluación}

El dise\~no contempla la evaluación contra etiquetas manuales. La infraestructura
almacena anotaciones humanas y calcula m\'etricas de precisi\'on, exhaustividad y
$F_1$, as\'\i{} como la $\kappa$ de Cohen para el acuerdo entre anotadores
\citep{cohen1960kappa,landis1977kappa}. En ausencia de etiquetas manuales, el informe
reporta las m\'etricas internas de cobertura y de nivel de evidencia, y se\~nala
explícitamente la necesidad de anotación futura.

\subsection{Desacoplamiento arquitectónico y servicio en Python/Flask}
\label{sec:flask_prod}

Aunque la experimentación, validación estadística y modelado temporal se desarrollaron en R (\texttt{engine.R}, prototipo interactivo en Shiny \texttt{app.R}), la necesidad de servir inferencias afectivas en tiempo real con alta concurrencia condujo al desacoplamiento de la lógica hacia un microservicio REST en \textbf{Python y Flask}. 

El servicio en Flask carga los recursos tabulares en memoria durante el arranque, exponiendo un endpoint parametrizado (\texttt{POST /analizar}) que recibe el texto y el código de idioma (\texttt{"lang": "es"|"en"}). Esta transición asegura tiempos de respuesta menores a 25 milisegundos por mensaje y permite integrar el clasificador determinista de 15 emociones en cualquier plataforma web moderna.